\documentclass[../../../include/open-logic-section]{subfiles}

\begin{document}

\olfileid[id]{his}{bio}{pet}

\olsection{R\'ozsa P\'eter}

\olphoto{peter-rozsa}{R\'ozsa P\'eter}

R\'ozsa P\'eter terlahir dengan nama R\'ozsa Politzer di Budapest,
Hungaria, pada 17 Februari 1905. Ia paling dikenal karena karyanya tentang
fungsi rekursif, yang sangat penting bagi lahirnya bidang teori rekursi.

P\'eter tumbuh dalam masa politik yang keras---Perang Dunia I berkecamuk
ketika ia remaja---tetapi dapat bersekolah di Sekolah Putri Maria Terezia
yang bergengsi di Budapest, tempat ia lulus pada tahun 1922. Ia kemudian
belajar di Universitas P\'azm\'any P\'eter (yang kelak berganti nama menjadi
Universitas Lor\'and E\"otv\"os) di Budapest. Atas desakan ayahnya, ia
mula-mula mempelajari kimia, tetapi kemudian beralih ke matematika dan lulus
pada tahun 1927. Meskipun memiliki kualifikasi untuk mengajar matematika
sekolah menengah, keadaan ekonomi ketika itu sangat buruk karena Depresi
Besar memengaruhi perekonomian dunia. Pada masa ini, P\'eter mengambil
berbagai pekerjaan lepas sebagai tutor dan guru privat matematika. Akhirnya
ia kembali ke universitas untuk menempuh studi pascasarjana dalam
matematika. Mula-mula ia berencana bekerja dalam teori bilangan, tetapi
setelah mengetahui bahwa hasil-hasilnya telah dibuktikan sebelumnya, ia
hampir meninggalkan matematika sama sekali. Ia didorong untuk meneliti
teorema ketaklengkapan G\"odel dan, tanpa menyadarinya, membuktikan beberapa
hasilnya dengan cara berbeda. Hal ini memulihkan kepercayaan dirinya,
dan P\'eter kemudian menulis makalah-makalah pertamanya tentang teori
rekursi, terilhami oleh program landasan David Hilbert. Ia memperoleh gelar
PhD pada tahun 1935, dan pada tahun 1937 menjadi editor \emph{Journal of
  Symbolic Logic}.

Makalah-makalah awal P\'eter secara luas dianggap sebagai sumbangan yang
meletakkan dasar bagi bidang teori fungsi rekursif. Dalam \cite{Peter1935a},
ia menyelidiki hubungan antara berbagai jenis rekursi. Dalam
\cite{Peter1935b}, ia menunjukkan bahwa suatu fungsi tertentu yang
didefinisikan secara rekursif bukanlah rekursif primitif. Hasil ini
menyederhanakan hasil terdahulu Wilhelm Ackermann. Fungsi P\'eter yang lebih
sederhana inilah yang kini sering disebut fungsi Ackermann---dan kadang,
secara lebih tepat, fungsi Ackermann--P\'eter. Ia menulis buku pertama
tentang teori fungsi rekursif \citep{Peter1951}.

Terlepas dari arti penting dan pengaruh karyanya, P\'eter baru memperoleh
posisi mengajar penuh waktu pada tahun 1945. Selama pendudukan Nazi di
Hungaria pada Perang Dunia II, P\'eter tidak diizinkan mengajar karena
undang-undang antisemit. Pada tahun 1944 pemerintah mendirikan sebuah ghetto
Yahudi di Budapest; ghetto itu dipisahkan dari bagian kota lainnya dan
dijaga oleh penjaga bersenjata. P\'eter dipaksa tinggal di ghetto tersebut
hingga dibebaskan pada tahun 1945. Ia kemudian mengajar di Kolese Pelatihan
Guru Budapest dan, sejak tahun 1955, di Universitas E\"otv\"os Lor\'and. Ia
adalah matematikawan perempuan Hungaria pertama yang menjadi Doktor Akademik
Matematika dan perempuan pertama yang terpilih menjadi anggota Akademi Ilmu
Pengetahuan Hungaria.

P\'{e}ter dikenal sebagai pengajar matematika yang penuh semangat, yang lebih
suka menjelajahi sifat dan keindahan masalah matematika bersama para
mahasiswanya daripada sekadar memberi kuliah. Karena itu, para mahasiswanya
dengan penuh kasih memanggilnya ``Bibi Rosa.'' P\'eter meninggal pada tahun
1977 dalam usia~71 tahun.

\begin{reading}
Untuk bacaan biografis lebih lanjut, lihat \citep{Oconnor2014} dan
\citep{Andrasfai1986}. \citet{Tamassy1994} melakukan wawancara singkat dengan
P\'eter. Untuk bacaan matematika yang menyenangkan, lihat buku P\'eter
\emph{Playing With Infinity} \citep{Peter2010}.
\end{reading}

\end{document}
